
\documentclass[11pt,oneside]{book}
\usepackage{../preamble}

\title{Idea Theory and the Social Sciences\\
\large Volume 2 of the IdeaTheory monograph}
\author{The IdeaTheory Project}
\date{\today}

\begin{document}
\frontmatter
\maketitle

\begin{abstract}
This is Volume~2 of the six-volume IdeaTheory monograph. Volume~1
established nine purely mathematical theorems about a graded monoid
$(\mathcal{I},\circ,e)$ of \emph{ideas} equipped with a bilinear
resonance pairing $R\colon\mathcal{I}\times\mathcal{I}\to\mathbb{R}$,
an Aufhebung decomposition $a\circ b = \mathrm{pres}(a,b) +
\mathrm{neg}(a,b) + \mathrm{lift}(a,b)$, an emergence $2$-cocycle
$\varepsilon$, a meaning-curvature form $\kappa$, conjugation, transmission
chains, structural equivalence, and a compatible grading. All nine
theorems were stated, proved, and machine-verified in Lean~4 in
Volume~1; the present volume contains \emph{no new mathematics}. Its
purpose is to interpret the formal apparatus across the social
sciences, so that every artifact in the IdeaTheory project visibly
bridges the existing non-formalized literature on the social life of
ideas to the machine-verified theorems.

The volume contains two chapters. Chapter~3, \emph{Ideas in Society:
Composition, Resonance, and Aufhebung}, reads idea composition $a\circ
b$ as the social combination of cultural units (memes, frames, norms,
institutional templates), the resonance pairing $R(a,b)$ as a
quantitative model of the alignment, tension, or borrowing between
ideas circulating in a population, and the Aufhebung decomposition as
a precise version of what historians of ideas, sociologists of
knowledge, and political theorists call \emph{preservation},
\emph{negation}, and \emph{sublation} when one tradition absorbs
another. We connect Theorem~1 (associativity of $\circ$) to coalition
algebra, Theorem~2 (bilinearity of $R$) to bibliometric and
citation-network models, and Theorem~3 (uniqueness of the Aufhebung
decomposition) to debates over hybridization and syncretism.
Chapter~4, \emph{Transmission, Conjugation, and Structural
Equivalence in Cultural and Political Life}, uses Theorems~6--8 to
formalize cultural transmission chains à la Sperber and Boyd, the
conjugation action $g\!\cdot\!a\!\cdot\!g^{-1}$ as translation between
discursive frames, and structural equivalence $a\sim b$ as the
sociological notion that two ideas play the same role in different
networks even when their surface content differs.

Each chapter opens with a \emph{Background from the Literature}
section citing the relevant non-formal sources (Durkheim, Weber,
Bourdieu, Sperber, Latour, Sewell, Boyd \& Richerson, and others),
followed by a \emph{From Informal Claims to Formal Statements} table
that maps each cited claim to a specific Lean theorem from
Volume~1. The volume closes with an appendix listing the Lean~4
files relevant to these interpretations and a paragraph for each.
No proofs are restated; the reader is referred to Volume~1 and to the
machine-checked Lean development.
\end{abstract}

\tableofcontents

\chapter*{Preface to Volume 2}
\addcontentsline{toc}{chapter}{Preface to Volume 2}

\section*{Motivation}
The mathematics of Volume~1 is austere. Nine theorems describe how
elements of an abstract graded monoid combine, pair, decompose,
deform, conjugate, and equilibrate. Read in isolation, that volume
could be mistaken for a fragment of noncommutative algebra with a
cohomological flavor. The IdeaTheory project, however, was conceived
as a bridge: a way to give the long, sprawling, and often imprecise
literature on \emph{ideas as social objects} a verifiable spine. The
present volume is the first half of that bridge, focused on the
social sciences. Its job is not to extend the formal theory but to
exhibit the social-scientific phenomena that motivate each axiom and
each theorem, and to record, claim by claim, how an existing piece of
sociological, anthropological, or political-theoretic discourse maps
onto a Lean-checked statement.

\section*{Intended audience}
The volume is written for two audiences in parallel. The first is the
\emph{social scientist} — sociologist of knowledge, cultural
anthropologist, political theorist, intellectual historian, or
science-and-technology-studies scholar — who is comfortable reading
formal definitions but does not want to read Lean source. For this
reader, every formal statement is paraphrased in prose and tied to a
recognizable debate in the literature. The second is the
\emph{formal-methods reader} who has worked through Volume~1 and
wishes to see why the axioms were chosen the way they were. For this
reader, each chapter contains a translation table whose right column
is a Lean lemma name and whose left column is a verbal claim drawn
from a cited source.

\section*{Prerequisites}
We assume familiarity with Volume~1 at the level of statements (not
proofs): the reader should know what is meant by the composition
$a\circ b$, the pairing $R(a,b)$, the Aufhebung triple
$(\mathrm{pres},\mathrm{neg},\mathrm{lift})$, the emergence cocycle
$\varepsilon$, and the curvature $\kappa$. From the social-science
side we assume only what an advanced undergraduate would know:
Durkheim's distinction between mechanical and organic solidarity,
Weber's ideal types, Bourdieu's notion of a field, the Sperber/Boyd
strand of cultural-evolution theory, and the basic vocabulary of
network sociology.

\section*{Reading order}
Volume~2 may be read immediately after Volume~1 or in parallel with
Volumes~3--6. Within Volume~2, Chapter~3 should be read before
Chapter~4 because the latter uses the Aufhebung decomposition
introduced in the former. Readers primarily interested in cultural
transmission may skim Chapter~3 §3.1--§3.2 and proceed directly to
Chapter~4 §4.2.

\section*{Relation to the other volumes}
Volume~1 supplies all the mathematics. Volume~3 takes up the same
formal apparatus for the humanities (literary, art-historical, and
hermeneutic interpretations). Volume~4 turns to cognitive science and
philosophy of mind, where ideas become mental representations and the
resonance pairing acquires a psycholinguistic reading. Volume~5
addresses the philosophical \emph{problem} of emergence using the
$2$-cocycle $\varepsilon$ in earnest. Volume~6 collects advanced and
unifying applications, including economic and computational
interpretations. The present volume is therefore the
social-scientific panel of a six-panel diptych whose mathematical
content is fixed once and for all in Volume~1.

\section*{Guided tour of the chapters}
Chapter~3 opens with a literature review of how sociologists and
historians of ideas have spoken about the combination of cultural
units, then formalizes the three classical operations
(\emph{preservation}, \emph{negation}, \emph{sublation}) as the
Aufhebung decomposition of Theorem~3, and shows that bilinearity of
$R$ (Theorem~2) is exactly what bibliometric and citation-network
models implicitly assume. Chapter~4 builds on this to read
transmission chains (Theorem~7), conjugation (Theorem~6), and
structural equivalence (Theorem~8) as formal counterparts of cultural
transmission, frame translation, and role equivalence in social
networks. Each chapter ends with an explicit
informal-to-formal translation table and a list of open questions
that the social-science literature poses to the formal theory.

\mainmatter

\input{ch3}
\input{ch4}

\appendix
\chapter{Lean 4 Files Relevant to Volume 2}

This appendix lists the Lean~4 source files in the IdeaTheory
formalization that are directly cited by the chapters of this volume.
All files live in the \texttt{IdeaTheory} namespace; full proofs are
in Volume~1 and in the machine-checked Lean development.

\paragraph{\texttt{IdeaTheory/Composition.lean}.}
Defines the graded monoid $(\mathcal{I},\circ,e)$ of ideas and proves
Theorem~1 (associativity, unit laws, grading compatibility). Cited in
Chapter~3 §3.1 as the formal counterpart of the social combination of
cultural units.

\paragraph{\texttt{IdeaTheory/Resonance.lean}.}
Defines the bilinear resonance pairing
$R\colon\mathcal{I}\times\mathcal{I}\to\mathbb{R}$ and proves
Theorem~2 (bilinearity and its consequences). Cited in Chapter~3 §3.2
as the algebraic backbone of bibliometric and citation-network
models.

\paragraph{\texttt{IdeaTheory/Aufhebung.lean}.}
Establishes the Aufhebung decomposition $a\circ b = \mathrm{pres}(a,b)
+ \mathrm{neg}(a,b) + \mathrm{lift}(a,b)$ and proves Theorem~3
(existence and uniqueness). Cited throughout Chapter~3 as the formal
version of preservation/negation/sublation in historical sociology.

\paragraph{\texttt{IdeaTheory/Emergence.lean}.}
Defines the emergence $2$-cocycle $\varepsilon$ and proves Theorem~4
(the cocycle identity). Cited in Chapter~3 §3.4 in connection with
collective effects in social aggregation; Volume~5 develops the
emergence reading in detail.

\paragraph{\texttt{IdeaTheory/Curvature.lean}.}
Introduces the meaning-curvature form $\kappa$ and proves Theorem~5
(asymmetry of $R$ measured by $\kappa$). Cited in Chapter~4 §4.1 as
the formal counterpart of asymmetric influence in directed
citation/discourse networks.

\paragraph{\texttt{IdeaTheory/Conjugation.lean}.}
Defines the conjugation action and proves Theorem~6 (its functorial
properties). Cited in Chapter~4 §4.2 as the formal model of frame
translation between discursive communities.

\paragraph{\texttt{IdeaTheory/Chains.lean}.}
Develops transmission chains and proves Theorem~7 (compositional
identity along chains). Cited in Chapter~4 §4.3 as the formal version
of cultural transmission in the Sperber/Boyd--Richerson tradition.

\paragraph{\texttt{IdeaTheory/Equivalence.lean}.}
Defines structural equivalence $a\sim b$ and proves Theorem~8 (it is
a congruence with respect to $\circ$ and $R$). Cited in Chapter~4
§4.4 as the formal version of role equivalence in network sociology.

\paragraph{\texttt{IdeaTheory/Grading.lean}.}
Establishes the grading and proves Theorem~9 (compatibility of all
the above structures with the grading). Cited in Chapter~4 §4.5 in
connection with stratified models of cultural complexity.

\bibliographystyle{alpha}
\bibliography{../references}

\end{document}
=== END FILE ===
